\chapter{Calibrating Coverage}
\label{ch:coverage}
The resolvent supplies a leading-order response to a small smooth surrogate
perturbation. It retains repeated linear feedback but leaves a prediction
error to measure. Calibration uses that remaining error to construct intervals
for a specified physical or numerical target.

\section{Why Conformal Calibration}
\label{sec:why-conformal}
Fix the predictor, observable, and nonconformity score before using calibration
labels. If calibration and test scores are exchangeable, split conformal provides
a finite-sample marginal coverage guarantee. It does not guarantee the same
coverage separately for every physical regime or for every individual state.
Changing the deployment population requires separate justification or recalibration.

\section{Residual Ensembles and Data Separation}
\label{sec:residual-ensemble}
For each calibration configuration, construct the corrected prediction using
that configuration's available surrogate equilibrium, residual, and derivative
callbacks. Independently solve the reference problem to obtain its target.
The input to \code{split\_conformal} is the remaining reference-minus-prediction
error, not simply the raw one-call surrogate residual.
Training a network or fitting PCA uses a separate training partition.
Hold calibration labels out of predictor selection and hyperparameter tuning.

\section{Physical Sampling}
\label{sec:stratified-sampling}
Declare the physical population and draw calibration and test configurations by
a compatible protocol. A positional split of an ordered parameter grid can
produce different regimes in the two partitions. Stratification may balance a
designed benchmark, but it does not by itself establish exchangeability or
coverage for a different production population. Retain failed configurations
and state how they affect the prediction rule and evaluation population.

\section{Calling Split Conformal}
\label{sec:calling-conformal}
\code{split\_conformal(calibration\_residuals, test\_predictions, alpha=...)}
returns a \code{ConformalBand} with \code{low}, \code{high},
\code{alpha}, and \code{n\_calibration}.
For vector targets the default score is the maximum absolute component error.
One scalar threshold then defines a simultaneous coordinate box.
The following complete example uses independent synthetic errors and fixed zero
predictions; its coverage is illustrative, not a physical benchmark:
\begin{lstlisting}[style=grailInStyle]
python examples/calibration.py
\end{lstlisting}
It uses 100 calibration vectors and 128 independent test vectors. The example
checks the threshold against a separately computed sorted order statistic and
also tests the infinite-band case for two calibration samples.

\section{Reading a Coverage Band}
\label{sec:reading-coverage}
The band surrounds the chosen prediction center. A GRAIL-corrected center is
$x_e-\delta x^{*,\mathrm{pred}}$. Narrower calibrated bands can reflect a
more accurate correction; they do not imply that physical feedback itself is
weak. Compare prediction methods using the same targets, score, sampling
population, and nominal coverage target. Actual coverage fractions fluctuate
in finite test sets.
The bands quantify the chosen surrogate-error target. Physical-model and
discretization uncertainties require their own treatment.

\section{PCA and Solve Count}
\label{sec:pca-coverage}
For one fixed linear response operator, propagate $r$ retained basis vectors
once and cache their responses. Propagate a nonzero residual mean separately,
then combine weights to approximate additional residual responses. This can
reduce the number of distinct linear solves from one per residual to $r$ plus
a mean solve. It is a solve-count comparison, not a measured wall-clock speedup.
\code{propagate} calls its supplied response function once per retained basis
vector on each invocation; the caller must arrange caching when reusing it.
Different physical configurations generally have different operators, so their
propagated bases cannot automatically share the same cache.
Fit the basis using training residuals only. Truncation error remains part of
the calibrated prediction error.

\section{Exact Finite-Sample Threshold}
\label{sec:finite-sample-limits}
For $N_{\rm cal}$ calibration scores and miscoverage $\alpha$, use
\begin{equation}
 k=\left\lceil(N_{\rm cal}+1)(1-\alpha)\right\rceil.
\end{equation}
The radius is the $k$th sorted score when $k\le N_{\rm cal}$ and infinity
otherwise. At nominal 90\% with two calibration samples, $k=3>2$, so the
valid band is infinite. The extra order-statistic correction does not imply
a universal percentage change in width; score spacing and tails determine
that change. Choose sample sizes for the required application and report the
observed width and held-out coverage.
